\documentclass[a4paper,10pt]{report}\input{../0.Préambule/Préambule_cours_prof}\begin{document}

\definecolor{titlecolor}{rgb}{0.12, 0.56, 1.0}
\newpage
\setcounter{chapter}{5}
\chapter{Calcul littéral}

\section{Simplification d'une expression littérale}

\begin{dft}
Une expression littérale est une expression mathématiques qui comporte une ou plusieurs lettres ainsi que des opérations. Ces lettres désignent des nombres.
\end{dft}

\subsection{Convention d'écriture avec le signe $\times$}

\begin{regle}
Pour simplifier l'écriture d'une expression littérale, on peut supprimer le symbole $\times$ devant une lettre ou une parenthèse.
\end{regle}

\begin{rmq}
On ne peut pas supprimer le signe $\times $ entre deux nombres.
\end{rmq}

\begin{ex}
Simplifie l'expression suivante : $A=-5 \times x+7 \times (-4) \times (3 \times x-2)$

\noindent\grandscarreaux{\linewidth}{2.4cm}
\end{ex}

\begin{regle}
Pour tout nombre $a$, on peut écrire : 
\begin{enumerate}
\item[•] $a \times a = a^2$ (qui se lit \og $a$ carré \fg{}) ;
\item[•] $a \times a \times a = a^3$(qui se lit \og $a$ au cube \fg{})
\end{enumerate}
\end{regle}

\subsection{Opposés et parenthèses}

\begin{prop}
L'opposé d'une somme algébrique est égal à la somme des opposés de chacun de ses termes.
\end{prop}

\begin{ex}
Quel est l'opposé de la somme algébrique $a+b-2ab$ ?

\noindent\grandscarreaux{\linewidth}{2.4cm}
\end{ex}

\begin{rmq}
Cette propriété permet de supprimer des parenthèses précédées d'un signe \og - \fg{} dans une expression.
\end{rmq}

\begin{ex}

\vspace{-5mm} \hspace{3cm} Supprime les parenthèses dans l'expression $B=3x-(-2x^2-5xy+4)$

\vspace{1mm}
\noindent\grandscarreaux{\linewidth}{2.4cm}
\end{ex}

\vspace{-3mm}
\section{Réduction d'une expression littérale}

\begin{dft}
Réduire une expression littérale, c'est écrire sous la forme d'une somme algébrique ayant le moins de termes possible.
\end{dft}

\begin{meth}[Réduction de somme]

\vspace{-3mm}
\noindent \begin{minipage}{5cm}
$$A = 3a + 3 + 5a - 1 - 2a + 4$$

\noindent\grandscarreaux{\linewidth}{3.2cm}
\end{minipage}
\hspace{5mm}
\begin{minipage}{5cm}
$$B = 3x^2 + 2x + 3 + x^2 + 5x + 7$$

\noindent\grandscarreaux{\linewidth}{3.2cm}
\end{minipage}
\hspace{5mm}
\begin{minipage}{5cm}
$$C = (x+8) - (6x-1) - (7-x)$$

\noindent\grandscarreaux{\linewidth}{3.2cm}
\end{minipage}
\end{meth}

\begin{meth}[Réduction de produit]

\noindent \begin{minipage}{5cm}
$$D = 3x \times 4$$

\noindent\grandscarreaux{\linewidth}{2.4cm}
\end{minipage}
\hspace{5mm}
\begin{minipage}{5cm}
$$E = 2a \times 3b$$

\noindent\grandscarreaux{\linewidth}{2.4cm}
\end{minipage}
\hspace{5mm}
\begin{minipage}{5cm}
$$F = 6y \times 3y$$

\noindent\grandscarreaux{\linewidth}{2.4cm}
\end{minipage}
\end{meth}

\section{Remplacer une lettre par une valeur dans une expression}

\begin{dft}
Calculer la valeur d'une expression littérale, c'est attribuer un nombre à chaque lettre de l'expression afin d'effectuer le calcul
\end{dft}

\begin{rmq}

\vspace{-4mm} \hspace{3cm} Une même lettre désigne toujours un même nombre dans une expression littérale donnée.
\end{rmq}

\begin{ex}

\vspace{-5mm} \hspace{3cm} Calculons \quad $A = 2 + 3t$ pour $t = 2$ \quad et  \quad $B=2x^2-3x+1$ pour $x=3$.

\vspace{1mm}
\noindent\grandscarreaux{\linewidth}{3.2cm}
\end{ex}


\end{document}